\begingroup

\section{Hochschild trace, Laplacian, and specialization: source integration}


\subsection{Hochschild trace, Laplacian, and specialization: source integration}
\label{sec:connes-pr27-source-integration}

The three complete notes adjoining this introduction are the
\emph{Hochschild comparison}, equations H1--H42, the
\emph{Laplacian parabola}, equations P1--P8, and the
\emph{Specialization research note}, equations S1--S28.
Their full definitions, proofs, source references, exceptions, and
verification qualifications are retained. They were reviewed at
commit \nolinkurl{a4a87844c257485ae6d669f5b3eb6d96b9e8b8d8}
and merged with the identical tree at
\nolinkurl{bf0ef873dfbfdbcba3bd0ec36f3442e403403eb8}.
The supplied earlier draft and checkpoint descriptions remain historical
source statements. The accompanying owner receipts record the accepted
reviewed-head checks and exact merged-source comparison.

The comparison uses Connes and Consani's
\emph{Hochschild homology, trace map and $\zeta$-cycles},
printed pages~83--101. The owner's complete source review read that
chapter and visually inspected printed pages~89--92, 94--95 and~100.
Its exact reading extent and the primary published-offprint hash are
retained in the evidence. No additional full-chapter inspection or
local formal execution is claimed by this integration.

\subsubsection{The precise trace image and the original source}

The factorization in H6 is read with its stated restricted target:
\[
 V\xrightarrow{j}j(V)\xrightarrow{\operatorname{Tr}}\mathscr B,
 \qquad \operatorname{Tr}j|_V=\Theta,
 \qquad \mathcal E=\tfrac12\mathcal U\Theta,
 \qquad(\mathcal UF)(u)=u^{1/2}F(u).
\]
The target $j(V)$ records the qualified trace image; it does not
assert an unrestricted strong-Schwartz-valued trace on every
Hochschild class. The isometry is precisely
$\mathcal U:L^2(du)\to L^2(du/u)$, whose norm equality is
$\int|u^{1/2}F(u)|^2du/u=\int|F(u)|^2du$.

The Gaussian $g_0=e^{-\pi u^2}$ belongs to the larger even Schwartz
space, and its image under $D^2-D$ is the original
$\phi_*=(4\pi^2u^4-6\pi u^2)e^{-\pi u^2}\in V$.
The full cochain comparison is H12, with both ordered legs in H13.
Its Fourier identity is exactly
$\mathcal FD=(1-D)\mathcal F$; physical parity instead commutes
with $D$. The original factor $\mathcal M\Theta\phi_*=2\xi$ remains.

Before applying the cited earlier continuous retraction, H15 retains
the map and kernel
\[
 \mathscr B/\Theta V\longrightarrow
 \mathscr B/\overline{\Theta V},\qquad
 \ker=\overline{\Theta V}/\Theta V.
\]
The retraction theorem is an earlier analytic input, separately
identified in the full note. A finite matrix calculation does not
prove its closed-image conclusion.

\subsubsection{The finite kernel and the retained zero mode}

The critical-line comparison H19 retains the plus-sign Mellin
convention, its action equations H20, the entire off-critical
generalized kernel H22, and the critical recovery map
\[
 \Xi_\rho[f]=\sum_{j<m}\frac{i^{-j}}{j!}
 f^{(j)}(\gamma)\epsilon^j,
 \qquad\rho=\tfrac12+i\gamma.
\]
The reciprocal multiplier proof applies to all nilpotent powers;
the surviving derivatives recover every critical jet.
Killed represented coefficients have value $(\lambda,0)$ in
the receiving support fibre. The map does not send them to
structural absence.

The raw zero coefficient $\xi(1/2)L^{-1/2}$ is nonzero and does
not extend smoothly to $L=0$. H34--H37 retain the extra smooth
scalar coordinate first, with reconstruction $L^{-1/2}a(L)$ in
the normalized Fourier basis and circle-function value $a(L)/L$.
The covering factor $\sqrt n$ gives the coordinate action
$a(L)\mapsto a(nL)$.

The subsequent original-source Gaussian average is H38. Its
large-$T$ convergence is in the stated local smooth-flat section
topology, with the full bound
\[
 p_{N,m,[0,M]}
 \bigl(\Pi_{\ne0}\widehat{\Sigma\mathcal E\phi_T}\bigr)
 \le C_{N,m,M}(1+T)^m e^{-2\pi^2T/M^2},\qquad T\ge1.
\]
The source does not assert uniform Schwartz seminorms for
$\phi_T$ as $T\to\infty$. The limiting pair
$(0,\xi(1/2))$, local smooth multiplication, and the closed-module
argument in H41 prove that the retained zero line becomes a
supported boundary in the resulting quotient.

The primary-source limit correction on printed page~100 must
remain explicit:
\[
 \zeta(1/2\pm2\pi i/L)\longrightarrow\zeta(1/2)\ne0
 \quad\text{as }L\longrightarrow\infty.
\]
The printed $L\to0$ at this step is erroneous. The owner checked
the actual rendered published page; the correction is present in
the complete detailed comparison note and is preserved here.

\subsubsection{The same arithmetic metric and its exceptional fibre}

H29--H31 and P1--P8 retain the original positive Gram $G$, action
$A$, and boundary form $W_k=A^*G+GA-kG$. The Laplacian
$A^2-kA$ is related to the source metrics through the exact
adjoint and energy identities, including all nilpotent jet terms.
The two-sided original-metric control yields the full numerical-range
parabola. Rank two alone does not make its relative error small.
The source's proposed growing-packet limit $\varepsilon_k/k\to0$
is explicitly unproved; it is not added to the established estimates.

The specialization note retains the source-coordinate metric
$O_j$, $C=\delta R^*O_j\delta R=G_i-G_j$, and the exact
cost/return transformation. Its all-order curvature coefficients,
telescoping identity, logarithmic integral and finite certificate
are all included, through S28. Their successful fixed-block stopping
degree is not asserted to be uniform over changing packets.

The local comparison $\alpha(w)=P_0+wP_1$, the action pole
$w^{-1}P_1AP_0$, and the derived-fibre line
$E_1\otimes\mathcal O(-1)|_\infty$ remain in the full source.
For the original cyclic quotient, the exact hull and quotient are
$F_A=(d)/(\chi)$ and $E/F_A\cong\mathbb C[S]/(d)$, with
$d=\gcd(\chi,p_1,\ldots,p_m)$ and $s_A=r-\deg d$.
Transporting the original metric retains
$\det G_A=|w|^{-2s_A}\det G_{\mathscr M}$ and the negative
boundary atom $-s_A\delta_\infty$ in its curvature current.

The accompanying evidence records the owner's written review,
actual finite executions and accepted exact-head CI. The analytic
maps, heat limit, coherent-sheaf consequences and numerical-range
proof remain identified as written mathematics; they are not
silently included in the selected Lean declarations.

Primary reference:
\href{https://doi.org/10.1090/pspum/105/01896}
{Connes and Consani, \emph{Hochschild homology, trace map and
$\zeta$-cycles}, Proceedings of Symposia in Pure Mathematics 105
(2023), pages 83--101}.


\endgroup
